\begin{afterword}
\summaryhead

The author recalls a theory he held at eighteen, that consciousness involved quantum gravity. He had followed the rules he was taught: the theory made a testable prediction, and Science would have let him spend ten years testing it as long as he admitted defeat afterward. What Science does not do, he argues, is say which theory is right before the experiment.

Some questions have large consequences but no cheap test now. His examples are macroscopic quantum superpositions, the untested parts of evolutionary psychology, and cryonics. For cryonics, the absence of revivals is only weak evidence against it, because no one expects revival to be possible with current technology. People who dismiss cryonics as ``not scientific'' misread science, whose verdict is only ``not proven,'' while 150,000 people die each day.

Where experiments cannot yet decide, he concludes, one must reason rationally and reach the answer before the evidence forces it.

\responsehead

The title says science fails on certain questions. On the essay's main case, it concedes the opposite: ``this is not a case where science actually gives the \emph{wrong answer}.'' Science says ``not proven'' about cryonics, and that is correct. What the essay objects to is people who read ``not proven'' as ``false,'' a misreading it could have corrected in a paragraph. The rest of its length does something else: it moves the reader from the fact that science does not decide to the author's own position.

It does this without making the case. The absence-of-evidence argument shows only that the lack of revivals is weak evidence against cryonics. It supplies no evidence for it. The positive case is declined (``I'm not going to go into those at the moment''), and its key premise, that freezing preserves the pattern of synapses, is replaced by a pointer to the author's physics posts. What fills the gap is pressure: 150,000 deaths a day, ``family and friends,'' an ambulance analogy that assumes the ambulance reaches a hospital, and a ``calculated bet'' with no calculation.

The examples are selected, not found. Every case where ``rationality has to take over from science'' is a position the author holds: many-worlds, evolutionary psychology, cryonics and, in the closing parenthesis, global catastrophic risk. The essay offers no case where reasoning ahead of the evidence went wrong, though it contains one. The diet theory it endorses on ``strangely varied anecdotal evidence'' did nothing for the author when he tried it, by his own report a year later.

The founding story does not show what it is used for. The young author's prediction was Penrose and Hameroff's. And by his own 2007 account, the rules he now blames ``helped me get out of the hole.'' One published physical objection to the microtubule idea, Tegmark's decoherence estimate of 2000, was a calculation made without waiting for an experiment: the kind of prior reasoning the essay says Science lacks.

The essay's one substantial point, that some hypotheses are supported only as parts of a larger productive theory, is Lakatos's from 1970. It arrives with the announcement that ``we'd need a new kind of verdict.'' The essay does not cite the literature it claims to replace.

In short: an essay that concedes science gives the right answer on its main example, then uses the concession to license belief in the author's causes without making the positive case for any of them.
\end{afterword}
